\documentclass[12pt]{ximera}
\title{Exam 3 Review}
\begin{document}
\begin{abstract}
Review for an Exam 3 on fair division and the start of apportionment: lone-divider and lone-chooser, sealed bids, markers, Hamilton's method and the Alabama paradox, and Jefferson's method.
\end{abstract}
\maketitle

\section*{Exam 3 Review}

Topics: lone-divider, lone-chooser, sealed bids, the method of markers, and
apportionment by Hamilton's and Jefferson's methods. Once you are comfortable with the
material, try these with a stopwatch: you should be able to finish them in under $60$
minutes.

\begin{problem}
Four burglars A, B, C, D divide their haul (fair division works for criminals too). D is
the lone divider and made four shares; the others' bids are:
\begin{center}
\begin{tabular}{r|ccc}
 & A & B & C\\\hline
Acceptable & $\set{s_2,s_4}$ & $\set{s_1,s_2}$ & $\set{s_2,s_3,s_4}$
\end{tabular}
\end{center}
How many different fair allocations of all four shares are there? $\answer{5}$

\begin{feedback}
D accepts anything, so only A, B, C constrain the choice.
\begin{itemize}
\item A $s_2$: B must take $s_1$; C takes $s_3$ or $s_4$ ($2$ ways).
\item A $s_4$, B $s_1$: C takes $s_2$ or $s_3$ ($2$ ways).
\item A $s_4$, B $s_2$: C must take $s_3$ ($1$ way).
\end{itemize}
D gets whatever is left each time.
\end{feedback}
\end{problem}

\begin{problem}
Three friends share a cake that is one quarter each buttercream, chocolate ganache,
vanilla, and strawberry.
\begin{itemize}
\item Alex loves chocolate and buttercream equally, but hates strawberry and vanilla.
\item Benji loves chocolate and strawberry equally, but hates buttercream and vanilla.
\item Catarina loves chocolate and vanilla equally, but hates buttercream and strawberry.
\end{itemize}
They use the lone-chooser method. (1) Catarina cuts two shares of equal value to her: one
with half the chocolate, half the vanilla, and all the strawberry, and the other with the
rest. (2) Benji picks the share he prefers. (3) Catarina and Benji each cut their share
into three sub-shares. (4) Alex picks one sub-share from each.

\begin{prompt}
\textbf{(a)} What is Catarina's first share (half chocolate, half vanilla, all strawberry)
worth to Benji? $\answer{3/4}$

\smallskip
So Benji picks it, and Catarina keeps the other: half chocolate, half vanilla, and all
the buttercream.
\begin{feedback}
To Benji, chocolate and strawberry are worth $\frac12$ each, so the share is worth
$\frac14+\frac12=\frac34$. The other share is worth only $\frac14$ to him.
\end{feedback}
\end{prompt}

\begin{prompt}
\textbf{(b)} Catarina's share is worth $\frac12$ to her, so each sub-share must be worth
$\frac16$. She cuts one sub-share of pure chocolate and one of pure vanilla. What
fraction of the whole chocolate goes in her pure-chocolate sub-share? $\answer{1/3}$
\begin{feedback}
Chocolate is worth $\frac12$ to her, so $\frac13$ of it is worth $\frac16$. The same goes
for vanilla. The third sub-share is the rest: $\frac16$ of the chocolate, $\frac16$ of
the vanilla, and all the buttercream, also worth $\frac1{12}+\frac1{12}=\frac16$.

Benji's share is worth $\frac34$ to him, so he makes three sub-shares of $\frac14$ each:
all his chocolate ($\frac12$ of the whole), half the strawberry, and the other half of
the strawberry with his vanilla.
\end{feedback}
\end{prompt}

\begin{prompt}
\textbf{(c)} Alex picks Catarina's third sub-share and Benji's chocolate sub-share. What
is Alex's total, as a fraction of the cake's value to Alex? $\answer{5/6}$
\begin{feedback}
Catarina's third sub-share: $\frac16$ of the chocolate ($\frac16\cdot\frac12=\frac1{12}$)
plus all the buttercream ($\frac12$), for $\frac7{12}$. Benji's chocolate sub-share is
half the chocolate: $\frac14$. Total: $\frac7{12}+\frac3{12}=\frac{10}{12}=\frac56$. A
very good deal for Alex!
\end{feedback}
\end{prompt}
\end{problem}

\begin{problem}
Alex, Bill, and Carla share a sandwich with equal parts meatball (M) and veggie (V).
Alex likes meatball twice as much as veggie, Bill hates veggie, and Carla likes both
equally.

\begin{prompt}
\textbf{(a)} Which methods are suitable?
\begin{selectAll}
\choice[correct]{Lone-divider}
\choice[correct]{Lone-chooser}
\choice{Sealed bids}
\choice{Method of markers}
\end{selectAll}
\begin{feedback}
A sandwich is a continuous asset, so a continuous method is needed.
\end{feedback}
\end{prompt}

\begin{prompt}
\textbf{(b)} \textbf{Lone-divider.} Alex divides into $s_1=\frac M2$, $s_2=\frac M2$,
$s_3=V$. Which shares are acceptable to Bill and to Carla?

\begin{selectAll}
\choice[correct]{Bill: $s_1$}
\choice[correct]{Bill: $s_2$}
\choice{Bill: $s_3$}
\choice{Carla: $s_1$}
\choice{Carla: $s_2$}
\choice[correct]{Carla: $s_3$}
\end{selectAll}
\begin{feedback}
To Alex, $M=\frac23$ and $V=\frac13$, so all three shares are worth $\frac13$. Bill values
each meatball half at $\frac12$ and the veggie at $0$. Carla values the meatball halves at
$\frac14$ each and the veggie at $\frac12$. So Carla gets $s_3$, Bill takes either
meatball half, and Alex gets the other.
\end{feedback}
\end{prompt}

\begin{prompt}
\textbf{(c)} \textbf{Lone-chooser}, with Alex as chooser. Bill cuts $\frac M2$ and
$\frac M2+V$; Carla takes $\frac M2+V$, worth $\frac34$ to her. Bill cuts his half into
three sixths of the meatball; Carla cuts hers into $\frac M2$, $\frac V2$, $\frac V2$.
Alex takes one of Bill's sixths and Carla's meatball half. What fraction of the meatball
does Alex end up with, and what is that worth to Alex?

Fraction of meatball: $\answer{2/3}$ \quad Worth to Alex: $\answer{4/9}$
\begin{feedback}
$\frac16+\frac12=\frac23$ of the meatball, worth $\frac23\cdot\frac23=\frac49$ to Alex.
Bill keeps $\frac13$ of the meatball, which is $\frac13$ in his view, and Carla keeps all
the veggie, worth $\frac12$ to her.
\end{feedback}
\end{prompt}
\end{problem}

\begin{problem}
Three people bid on the items at an estate sale, using sealed bids.
\begin{center}
\begin{tabular}{r|rrr}
 & A & B & C\\\hline
Cabinet & 100 & 100 & 200\\
Table & 200 & 100 & 100\\
Bookshelf & 50 & 100 & 50
\end{tabular}
\end{center}

\begin{prompt}
\textbf{(a)} Who gets each item?

Cabinet $\answer{C}$ \quad Table $\answer{A}$ \quad Bookshelf $\answer{B}$
\end{prompt}

\begin{prompt}
\textbf{(b)} Initial settlement (positive = pays), to the nearest cent:

A $\answer[tolerance=0.005]{83.33}$ \quad B $\answer{0}$ \quad C
$\answer[tolerance=0.005]{83.33}$
\begin{feedback}
Totals: A $350$, B $300$, C $350$, so fair shares are $116.67$, $100$, and $116.67$. A and C
each receive items worth $200$ and pay the $83.33$ excess; B's bookshelf is worth exactly
B's fair share.
\end{feedback}
\end{prompt}

\begin{prompt}
\textbf{(c)} Each person's share of the surplus, to the nearest cent:
$\$\answer[tolerance=0.005]{55.56}$
\begin{feedback}
The surplus is $\frac{500}3\approx166.67$; a third of that is $\frac{500}9\approx55.56$.
\end{feedback}
\end{prompt}
\end{problem}

\begin{problem}
Housemates A, B, C use sealed bids to divide yardwork. The amounts are what each would
charge for the chore.
\begin{center}
\begin{tabular}{r|rrr}
 & A & B & C\\\hline
Leaves & 50 & 55 & 20\\
Grass & 40 & 15 & 30\\
Hedge & 30 & 35 & 40
\end{tabular}
\end{center}

\begin{prompt}
\textbf{(a)} Who does each chore?

Leaves $\answer{C}$ \quad Grass $\answer{B}$ \quad Hedge $\answer{A}$
\end{prompt}

\begin{prompt}
\textbf{(b)} Initial settlement (positive = pays): A $\answer{10}$ \quad B $\answer{20}$
\quad C $\answer{10}$
\begin{feedback}
Fair shares of the work are $40$, $35$, and $30$. A does $30$ worth, B $15$, and C $20$.
Everyone does less than a fair share in their own view, so everyone pays in.
\end{feedback}
\end{prompt}

\begin{prompt}
\textbf{(c)} Surplus: $\$\answer{40}$
\begin{feedback}
$10+20+10=40$, or $\$13.33$ each.
\end{feedback}
\end{prompt}
\end{problem}

\begin{problem}
There are $4$ Brach's (B), $5$ candy corn (C), and $6$ Droobles (D), lined up as follows:
\begin{center}
\begin{tabular}{ccccccccccccccc}
1 & 2 & 3 & 4 & 5 & 6 & 7 & 8 & 9 & 10 & 11 & 12 & 13 & 14 & 15\\\hline
B & C & D & D & D & B & C & D & D & B & B & C & C & D & C
\end{tabular}
\end{center}
Xavier likes Brach's best, Yvonne candy corn, and Zeynab Droobles, but each also wants as
many pieces as possible. They use the method of markers with two rules: (1) each share
holds between one-third of the favorite candy rounded down and one-third rounded up; (2)
subject to rule 1, shares are as balanced in piece count as possible.

\begin{prompt}
\textbf{(a)} Zeynab's markers are forced. After which pieces does she place them?

$\answer{4}$ and $\answer{8}$
\begin{feedback}
Six Droobles means exactly $2$ per share. The first share must contain pieces $3$ and $4$
but not $5$; the second must contain $5$ and $8$ but not $9$.
\end{feedback}
\end{prompt}

\begin{prompt}
\textbf{(b)} Xavier marks after $5$ and $10$, and Yvonne after $6$ and $12$. Complete the
allocation.

Zeynab gets $1$--$\answer{4}$ \quad Xavier gets $\answer{6}$--$\answer{10}$ \quad
Yvonne gets $\answer{13}$--$15$
\begin{feedback}
Zeynab's first marker is leftmost, so she takes $1$--$4$. Xavier's second marker
(after $10$) precedes Yvonne's (after $12$), so he takes his second share, $6$--$10$.
Yvonne takes her last share, $13$--$15$.
\end{feedback}
\end{prompt}

\begin{prompt}
\textbf{(c)} Which pieces are the surplus?
\begin{selectAll}
\choice[correct]{5}
\choice{6}
\choice{10}
\choice[correct]{11}
\choice[correct]{12}
\choice{13}
\end{selectAll}
\end{prompt}
\end{problem}

\begin{problem}
You are one of three players sharing assets worth $1,1,2,2,3,5,5,5$ in your view, lined
up in that order.

\begin{prompt}
\textbf{(a)} How much should each fair share be worth? $\answer{8}$
\end{prompt}

\begin{prompt}
\textbf{(b)} Why is the method of markers a problem in this order?
\begin{freeResponse}
\end{freeResponse}
\begin{feedback}
The running totals are $1,2,4,6,9,\ldots$, which skip $8$, and every later share is worth
at least $5$. There is no way to cut three consecutive shares worth $8$: one share is
always too small (as little as $5$) or too large (at least $10$).
\end{feedback}
\end{prompt}

\begin{prompt}
\textbf{(c)} Describe a rearrangement that allows three perfect shares, and what could
still go wrong.
\begin{freeResponse}
\end{freeResponse}
\begin{feedback}
For example, $1,2,5\mid1,2,5\mid3,5$. Still, the other players may value the items
differently, and there is no guarantee that everyone can find three acceptable shares in
the same order.
\end{feedback}
\end{prompt}
\end{problem}

\begin{problem}
The four states of Alphabitia use Hamilton's method to apportion the $60$ seats of their
General Assembly.
\begin{center}
\begin{tabular}{r|rrrrr}
 & A & B & C & D & Total\\\hline
Population & 152 & 18 & 92 & 38 & 300
\end{tabular}
\end{center}

\begin{prompt}
\textbf{(a)} Standard divisor $\answer{5}$; quotas A $\answer{30.4}$, B $\answer{3.6}$,
C $\answer{18.4}$, D $\answer{7.6}$
\end{prompt}

\begin{prompt}
\textbf{(b)} Hamilton apportionment: A $\answer{30}$, B $\answer{4}$, C $\answer{18}$,
D $\answer{8}$
\begin{feedback}
Lower quotas total $58$; the residues $0.6$ of B and D are the largest.
\end{feedback}
\end{prompt}

\begin{prompt}
\textbf{(c)} The Assembly grows to $61$ seats. Find the new standard quota for A and the
new Hamilton apportionment.

$q_A=\answer[tolerance=0.005]{30.91}$ \qquad A $\answer{31}$, B $\answer{3}$,
C $\answer{19}$, D $\answer{8}$
\begin{feedback}
$SD=\frac{300}{61}\approx4.918$, and the quotas are $30.91$, $3.66$, $18.71$, $7.73$. The
lower quotas total $58$, so $3$ seats go to the largest residues: A ($0.91$), D ($0.73$),
and C ($0.71$).
\end{feedback}
\end{prompt}

\begin{prompt}
\textbf{(d)} What happened?
\begin{multipleChoice}
\choice[correct]{The Alabama paradox: B lost a seat when the house grew}
\choice{The population paradox}
\choice{The new-states paradox}
\choice{A quota-rule violation}
\end{multipleChoice}
\end{prompt}
\end{problem}

\begin{problem}
Three provinces --- Alberta, Beringia, and Carolina --- form a new federation and use
Jefferson's method to apportion $50$ seats on their council.
\begin{center}
\begin{tabular}{r|rrr}
Divisor & Alberta & Beringia & Carolina\\\hline
Population & 9,470 & 3,899 & 6,631\\
$400$ & 23.675 & 9.748 & 16.578\\
$394$ & 24.04 & 9.90 & 16.83\\
$390$ & 24.282 & 9.997 & 17.003\\
$385$ & 24.60 & 10.13 & 17.22
\end{tabular}
\end{center}

\begin{prompt}
\textbf{(a)} Which divisor works? $\answer{390}$ \qquad Apportionment: Alberta
$\answer{24}$, Beringia $\answer{9}$, Carolina $\answer{17}$
\begin{feedback}
Rounding down: $400\to48$, $394\to49$, $390\to50$, $385\to51$.
\end{feedback}
\end{prompt}

\begin{prompt}
\textbf{(b)} What would have gone wrong had the $390$ row been rounded to two decimals
like the others?
\begin{freeResponse}
\end{freeResponse}
\begin{feedback}
$9.997$ would appear as $10.00$, and the lower quotas would seem to be $24+10+17=51$. We
would wrongly reject $390$. Whenever a rounded quota looks like an exact integer, use
more digits to see which side of the integer it is really on.
\end{feedback}
\end{prompt}
\end{problem}

\end{document}
